\documentclass[11pt,letterpaper]{article}
\usepackage{technical-user-manual}
\setDocTitle{Chapter 8: Continuous Deployment}
\setDocSubtitle{Foundation study guide | printed chapter page 189}
\setDocVersion{1.0}
\setDocStatus{Technical education / proposed reference}
\setDocOwner{CI/CD Engineering}
\setDocDate{30 October 2026}
\setDocSummary{Book-grounded original study guide; Chapter 8 plus explicitly marked extensions}
\begin{document}
\makedoctitle
\makedocsummary
\tableofcontents
\clearpage
\section{Reading identity and provenance}
Primary text: Paul M. Duvall with Steve Matyas and Andrew Glover, Continuous Integration: Improving Software Quality and Reducing Risk (2007), Chapter 8, Continuous Deployment; opening printed page 189. This study guide is an original pedagogical synthesis and does not reproduce chapter text.

Book-grounded concepts are presented in the next section. The workshop, cloud control implementations, modern terminology and assessment procedures are supplementary engineering material.

\section{Conceptual synthesis}
1. A releasable build requires traceable labels, a clean target environment, test evidence, distribution and recovery capability.

2. Rebuilding artifacts between environments can weaken traceability; promote immutable, integrity-checked artifacts.

3. Modern distinction: continuous delivery keeps release available; continuous deployment automatically releases eligible changes.

\section{System analysis and operational implications}
The central engineering question is how the Chapter 8 practices prevent changes from accumulating undetected integration risk. An effective answer names the triggering change, the evidence produced, the point of developer feedback and the recovery step.

Analyze four failure modes: implementation defect, test infrastructure defect, external dependency outage and policy violation. For each, specify an independent confirming observation and avoid declaring a release ready solely because a single status check is green.

Separate evidence generation from the release decision. Checks can supply objective data, but an accountable approver or explicit policy must decide the impact and permitted response.

\section{Guided exercise}
Complete the deployment architecture specification and demonstrate a simulated failed release decision.

\begin{enumerate}
\item State the initial system and versioned inputs.
\item Predict the expected failure or success signal before running a tool.
\item Execute locally, then in a disposable CI environment where relevant.
\item Capture command, revision, run reference, result and troubleshooting observations.
\item Explain one improvement that shortens detection without weakening assurance.
\end{enumerate}
\section{Knowledge check}
\begin{itemize}
\item Draw immutable artifact promotion from CI to protected staging and production.
\item Write a rollback decision tree with database compatibility constraints.
\item Identify which release actions require human approval or separation of duties.
\item Which book principle in Chapter 8 remains independent of a specific CI vendor?
\item Which contemporary tool or security control would you add today, and why is it an extension rather than a quotation from the 2007 text?
\end{itemize}
\section{Completion evidence and acceptance criteria}
\begin{itemize}
\item A timestamped experiment log identifies inputs and result.
\item A failure path and recovery action are described, not only the happy path.
\item The artifact links to its source revision and pipeline/workstation run.
\item At least one residual risk is stated with a responsible owner.
\item A peer can reproduce the exercise from the written instructions.
\end{itemize}
\end{document}
